
\documentclass[a4paper,11pt]{article}
\usepackage[a4paper, margin=8em]{geometry}

% usa i pacchetti per la scrittura in italiano
\usepackage[french,italian]{babel}
\usepackage[T1]{fontenc}
\usepackage[utf8]{inputenc}
\frenchspacing 

% usa i pacchetti per la formattazione matematica
\usepackage{amsmath, amssymb, amsthm, amsfonts}

% usa altri pacchetti
\usepackage{gensymb}
\usepackage{hyperref}
\usepackage{standalone}

% imposta il titolo
\title{Appunti Comunicazioni Numeriche}
\author{Luca Seggiani}
\date{2026}

% imposta lo stile
% usa helvetica
\usepackage[scaled]{helvet}
% usa palatino
\usepackage{palatino}
% usa un font monospazio guardabile
\usepackage{lmodern}

\renewcommand{\rmdefault}{ppl}
\renewcommand{\sfdefault}{phv}
\renewcommand{\ttdefault}{lmtt}

% disponi il titolo
\makeatletter
\renewcommand{\maketitle} {
	\begin{center} 
		\begin{minipage}[t]{.8\textwidth}
			\textsf{\huge\bfseries \@title} 
		\end{minipage}%
		\begin{minipage}[t]{.2\textwidth}
			\raggedleft \vspace{-1.65em}
			\textsf{\small \@author} \vfill
			\textsf{\small \@date}
		\end{minipage}
		\par
	\end{center}

	\thispagestyle{empty}
	\pagestyle{fancy}
}
\makeatother

% disponi teoremi
\usepackage{tcolorbox}
\newtcolorbox[auto counter, number within=section]{theorem}[2][]{%
	colback=blue!10, 
	colframe=blue!40!black, 
	sharp corners=northwest,
	fonttitle=\sffamily\bfseries, 
	title=~\thetcbcounter: #2, 
	#1
}

% disponi definizioni
\newtcolorbox[auto counter, number within=section]{definition}[2][]{%
	colback=red!10,
	colframe=red!40!black,
	sharp corners=northwest,
	fonttitle=\sffamily\bfseries,
	title=~\thetcbcounter: #2,
	#1
}

% U.D.M
\newcommand{\amp}{\ensuremath{\mathrm{A}}}
\newcommand{\volt}{\ensuremath{\mathrm{V}}}
\newcommand{\meter}{\ensuremath{\mathrm{m}}}
\newcommand{\second}{\ensuremath{\mathrm{s}}}
\newcommand{\farad}{\ensuremath{\mathrm{F}}}
\newcommand{\henry}{\ensuremath{\mathrm{H}}}
\newcommand{\siemens}{\ensuremath{\mathrm{S}}}

% circuiti
\usepackage{circuitikz}
\usetikzlibrary{babel}

% disegni
\usepackage{pgfplots}
\pgfplotsset{width=10cm,compat=1.9}

% disponi codice
\usepackage{listings}
\usepackage[table]{xcolor}

\lstdefinestyle{codestyle}{
		backgroundcolor=\color{black!5}, 
		commentstyle=\color{codegreen},
		keywordstyle=\bfseries\color{magenta},
		numberstyle=\sffamily\tiny\color{black!60},
		stringstyle=\color{green!50!black},
		basicstyle=\ttfamily\footnotesize,
		breakatwhitespace=false,         
		breaklines=true,                 
		captionpos=b,                    
		keepspaces=true,                 
		numbers=left,                    
		numbersep=5pt,                  
		showspaces=false,                
		showstringspaces=false,
		showtabs=false,                  
		tabsize=2
}

\lstdefinestyle{shellstyle}{
		backgroundcolor=\color{black!5}, 
		basicstyle=\ttfamily\footnotesize\color{black}, 
		commentstyle=\color{black}, 
		keywordstyle=\color{black},
		numberstyle=\color{black!5},
		stringstyle=\color{black}, 
		showspaces=false,
		showstringspaces=false, 
		showtabs=false, 
		tabsize=2, 
		numbers=none, 
		breaklines=true
}

\lstdefinelanguage{javascript}{
	keywords={typeof, new, true, false, catch, function, return, null, catch, switch, var, if, in, while, do, else, case, break},
	keywordstyle=\color{blue}\bfseries,
	ndkeywords={class, export, boolean, throw, implements, import, this},
	ndkeywordstyle=\color{darkgray}\bfseries,
	identifierstyle=\color{black},
	sensitive=false,
	comment=[l]{//},
	morecomment=[s]{/*}{*/},
	commentstyle=\color{purple}\ttfamily,
	stringstyle=\color{red}\ttfamily,
	morestring=[b]',
	morestring=[b]"
}

% disponi sezioni
\usepackage{titlesec}

\titleformat{\section}
	{\sffamily\Large\bfseries} 
	{\thesection}{1em}{} 
\titleformat{\subsection}
	{\sffamily\large\bfseries}   
	{\thesubsection}{1em}{} 
\titleformat{\subsubsection}
	{\sffamily\normalsize\bfseries} 
	{\thesubsubsection}{1em}{}

% disponi alberi
\usepackage{forest}

\forestset{
	rectstyle/.style={
		for tree={rectangle,draw,font=\large\sffamily}
	},
	roundstyle/.style={
		for tree={circle,draw,font=\large}
	}
}

% disponi algoritmi
\usepackage{algorithm}
\usepackage{algorithmic}
\makeatletter
\renewcommand{\ALG@name}{Algoritmo}
\makeatother

% disponi numeri di pagina
\usepackage{fancyhdr}
\fancyhf{} 
\fancyfoot[L]{\sffamily{\thepage}}

\makeatletter
\fancyhead[L]{\raisebox{1ex}[0pt][0pt]{\sffamily{\@title \ \@date}}} 
\fancyhead[R]{\raisebox{1ex}[0pt][0pt]{\sffamily{\@author}}}
\makeatother

\begin{document}
% sezione (data)
\section{Lezione del 14-05-26}

% stili pagina
\thispagestyle{empty}
\pagestyle{fancy}

% testo
\subsection{Processi di rumore bianco}
Approfondiamo il discorso del rumore rifacendoci ai \textit{processi aleatori}. Il \textbf{rumore bianco} a tempo continuo è un processo aleatorio che ha densità spettrale ugualmente distribuita sullo spettro frequenziale.

Definiamo il rumore bianco prendendo la funzione di autocorrelazione $R_X(\tau)$ triangolare, di base $\tau_{cor}$ e area $\xi$:
$$
R_X(\tau)
= \frac{\xi}{\tau_{cor}} \cdot \left( 1 - \frac{|\tau|}{\tau_{cor}} \right) \cdot \mathrm{rect} \left( \frac{|\tau|}{2\tau_{cor}} \right)
= \frac{\xi}{\tau_{cor}} \cdot \mathrm{tri} \left( \frac{\tau}{\tau_{cor}} \right)
$$
Dalla funzione di autocorrelazione ricordiamo possiamo ricavare la densità spettrale di potenza $S_X(f)$:
$$
S_X(f) = \xi \cdot \mathrm{sinc}^2 \left( f \cdot \tau_{cor} \right)
$$

Ora, facendo tendere $\tau_{cor} \rightarrow 0$ otteniamo che la funzione di autocorrelazione approssima sempre di più la delta di Dirac, mentre la densità spettrale di potenza approssima sempre di più la funzione costante con valore $\xi$. Si ha quindi:
$$
R_X(\tau) \rightarrow \xi \delta(\tau), \quad S_X(f) \rightarrow S_X(0) = \xi
$$

Da quanto avevamo detto riguardo alla funzione di autocorrelazione (e quindi la densità spettrale di potenza), abbiamo che:
\begin{itemize}
	\item Autocorrelazione più vicina a $0$ vicino all'origine rappresenta un processo che dà sempre minore informazione riguardo al suo comportamento locale; 
	\item Equivalentemente, densità spettrale di potenza più ampia rappresenta un processo che varia più velocemente (quindi ha maggiore banda); 
\end{itemize}
Il limite porta esattamente ad una autocorrelazione nulla ovunque tranne che in zero (la $\delta$), e una densità spettrale costante su tutte le frequenze: facendo questo limite prendiamo quindi una funzione (chiaramente idaele) che non dà alcuna informazione riguardo al suo comportamento passato, e ha banda infinita.

\subsection{Processi aleatori gaussiani a tempo continuo}
Di particolare interesse ci sono i processi dove le variabili aleatorie $[X(t_1), ..., X(t_n)]$ che compongono ogni insieme di istanti temporali $[t_1, ..., t_n]$ del processo sono \textit{congiuntamente gaussiane}.

Più formalmente:
\begin{definition}{Processo aleatorio gaussiano}
Un processo aleatorio viene detto di tipo gaussiano se, preso ogni insieme di istanti temporali:
$$[t_1, ..., t_n]$$
le variabili aleatorie:
$$[X(t_1), ..., X(t_n)]$$
da questo ricavato sono congiuntamente gaussiane.
\end{definition}

\begin{itemize}
\item Ricordiamo che se più variabili sono congiuntamente gaussiane, allora ognuna di esse è gaussiana;
\item Ricordiamo inoltre che se le variaili aleatorie gaussiane $[X(t_1), ..., X(t_n)]$ sono incorrelate (hanno covarianza nulla), allora sono anche indipendenti. 
\item Infine, una proprietà importante della densità di probabilità congiunta gaussiana è che se un processo gaussiano $X(t)$ è stazionario in senso lato, lo è automaticamente anche in senso stretto.
\end{itemize}

\subsubsection{Filtraggio di processi aleatori gaussiani}
Avevamo visto che il problema del filtraggio di un processo aleatorio non è completamente risolubile. Quindi, in generale è impossibile ottenere la descrizione completa del processo d’uscita $Y(t)$ nota quella del processo d’ingresso $X(t)$. I processi aleatori gaussiani sono l'eccezione a questa regola, in quanto filtrando un processo aleatorio gaussiano $X(t)$ si ottiene un secondo processo aleatorio gaussiano $Y(t)$.

Avevamo già calcolato la variazione degli indicatori relativi alla stazionarietà dei processi aleatori filtrati, per cui possiamo dire:
\begin{enumerate}
\item Prendendo un processo aleatorio gaussiano come stazionario, fissiamo i parametri della stazionarietà in senso lato (valor medio e autocorrelazione):
$$
\begin{cases}
\eta_x (t) = \eta_x \\
R_x(t_1, t_2) = R_x(\tau)
\end{cases}
$$
da cui, visto che $Y(t)$ è gaussiano:
$$
\implies \begin{cases}
\eta_x(t) = \eta_x(t) \\
\sigma_x^2(t_0) = R_x(t_0, t_0) - \eta_x^2(t)
\end{cases}, \quad
X(t_0) \in \mathrm{Gaussiana}\left( \eta_x, \sigma_x^2 \right)
$$
\item Invece, ponendo il filtraggio attraverso un dato filtro da risposta in frequenza $H$, abbiamo:
$$
\begin{cases}
\eta_y(t) = \eta_x(t) \circledast h(t) \\
R_y(t_0, t_1) = R_x(t_0, t_1) \circledast h(t_0) \circledast h(t_1) 
\end{cases}
$$
da cui, visto che $Y(t)$ è gaussiano:
$$
\implies \begin{cases}
\eta_y(t) = \eta_x(t) \circledast h(t) \\
\sigma_x^2(t_0) = R_y(t_0) - \eta_y^2(t_0) 
\end{cases}, \quad
Y(t_0) \in \mathrm{Gaussiana}\left( \eta_y, \sigma_y^2 \right)
$$
\item Prendendo un processo aleatorio in ingresso al filtro che è sia gaussiano che stazionario, siamo nella situazione più semplice di tutte, dove:
$$
\begin{cases}
\eta_y = \eta_x H(0) \\
R_y(\tau) = R_x(\tau) \circledast h(-\tau) \circledast h(\tau) 
\end{cases}
$$
da cui, visto che $Y(t)$ è gaussiano:
$$
\implies \begin{cases}
\eta_y(t) =  \eta_x(t) H(0) \\
\sigma_x^2(t) = R_y(0) - \eta_y^2(t)
\end{cases}, \quad
Y(t_0) \in \mathrm{Gaussiana}\left( \eta_y, \sigma_y^2 \right)
$$
\end{enumerate}

\subsubsection{Processo aleatorio gaussiano bianco}
Il processo aleatorio gaussiano è anche un buon modello per il processo rumore bianco, in quanto si può porre:
$$
\begin{cases}
	\eta_w = 0 \\
	\sigma_w^2 \rightarrow + \infty
\end{cases}
$$
ed ottenere effettivamente un processo di rimore bianco. Il parametro $\eta_w = 0$ è banale, mentre $\sigma_w^2$ si ricava ponendo:
$$
P_w = \sigma_w^2 + \eta_w^2 = \sigma_w^2 = \int_{-\infty}^{\infty} S_w(f) \, df = +\infty
$$
visto che dalla definizione di rumore bianco:
$$
R_w(\tau) = \xi \delta(\tau) , \quad S_w(f) = \xi
$$

\end{document}
